% Replace def:hb in full, including its unnumbered comparison paragraph.
% Place after night:P1.15, before the merged lifted-period theorem.
% F.18, F.19 and F.20 move here from the merger, so every label occurs once.
\begin{definition}[The normalized height function]
\label{def:hb}
Let $e=e_0$ be the fixed real section of the rank-six system,
and let $\sigma_b:P\to V$ be its algebraic splitting.
The invariant bivector, primitive projection and contraction are
normalized in the companion marking as in \cref{night:P1.15}.
For each regular base point, including the finite root points,
let $n$ be the real vector in $P\cap\bar P$ with $\alpha(n)=1$.
We define $h_b$ by
\begin{equation}
 \sigma_b(n)-\overline{\sigma_b(n)}=2i h_b e.
 \hypertarget{tag-02ER}{}\label{maps:eq:F.20}
\end{equation}
\end{definition}
\begin{proof}[Well-definedness]
Let $u$ be a local nonzero section of $\ell_0$.
Its Hodge metric is positive, so $\langle u,\bar u\rangle$
is a nondegenerate symplectic plane. The inverse polarization
bivector $\omega$ has a nonzero component on the exterior square
of its symplectic orthogonal complement. Consequently an equality
$u\wedge v-\bar u\wedge w=t\omega$ forces $t=0$; with $t=0$
its common value belongs to $\mathbb C(u\wedge\bar u)$.
The primitive projection therefore gives
\begin{equation}
 P\cap\bar P=p(\ell_0\wedge\bar\ell_0).
 \hypertarget{tag-02EP}{}\label{maps:eq:F.19}
\end{equation}
We make the scalar in the primitive decomposition explicit.
In the algebraic de Rham frame put
$\omega_{\rm alg}=3\omega_{03}+\omega_{12}$ and let
$\alpha_{\rm alg}(v)$ be the contraction of its unique wedge
lift $j(v)$ with the residue pairing $J_{03}=4/3$, $J_{12}=4$. Then
$\operatorname{contr}_{J}(\omega_{\rm alg})=8$. The raw primitive
coefficient is therefore $\alpha_0=\alpha_{\rm alg}/8$, whereas
\cref{night:P1.15} fixes
\[
 \alpha=\frac{625\pi i}{128}\alpha_{\rm alg}
        =\frac{625\pi i}{16}\alpha_0,
 \qquad
 \omega_{\rm norm}=\frac{16}{625\pi i}\omega_{\rm alg}.
\]
For the primitive representative $v$ of a quotient vector, its
unique wedge lift is exactly
\begin{equation}
 j(v)=v+\alpha_0(v)\omega_{\rm alg}
     =v+\alpha(v)\omega_{\rm norm}.
 \hypertarget{tag-02EO}{}\label{maps:eq:F.18}
\end{equation}

Here $\omega_{\rm norm}$ is real in the marked flat system,
although the algebraic formula contains $i$. To see the precise
comparison, use the companion cup matrix $\varepsilon J_c$ of
\cref{night:P1.7}, with $\varepsilon\in\{1,-1\}$, and its
invariant bivector $\omega_c=(1,1,1,1,1,1)^t$.
Contraction with $J_c$ gives $\operatorname{contr}_{J_c}(\omega_c)=2$.
For the rank-four period comparison $\Pi$, the residue formula is
$\Pi^t(\varepsilon J_c)\Pi=2\pi iJ$.
Indeed, cut the curve along a symplectic polygon and small circles
around the poles; the boundary integral of a primitive times the
second differential is the alternating period pairing on the
polygon and $2\pi i$ times the sum of residues on the circles.
It follows by inversion of the two alternating forms that
$(\bigwedge^2\Pi)\omega_{\rm alg}=(8\pi i/\varepsilon)\omega_c$.

The scaled exterior-square map
$\Psi=(625\varepsilon/256)\bigwedge^2\Pi$ induces the
quotient comparison in the fixed rank-six extension. Its scalar is forced by the cusp contraction identity: on the
kernel vector of \cref{night:P1.15} the extension residue is
$625/192$, while $\alpha_{\rm alg}$ has value $4/3$; the
logarithmic monodromy contributes $2\pi i$. The ratio is
$625\pi i/128$, exactly the pullback of
$\operatorname{contr}_{J_c}$ under this scaled comparison.
The one-dimensional extension calculation in
\cref{night:P1.15} then lifts it uniquely to the normalized
rank-six comparison fixing $e$.
Consequently $\Psi(\omega_{\rm norm})=\omega_c/2$, and
the induced primitive comparison sends $\alpha$ to
$\operatorname{contr}_{J_c}$. Both are real in the companion
marking. For $r=ip(u\wedge\bar u)$, the value $\alpha(r)$
is real and nonzero by the positive polarization pairing.
Thus $n=r/\alpha(r)$ is the unique real vector specified above.
Multiplying $u$ by a nonzero complex number multiplies both
$r$ and $\alpha(r)$ by its squared modulus, and leaves $n$
unchanged.

The image of $\sigma_b(n)$ in the quotient is real. Its
difference from its conjugate is therefore an imaginary
multiple of the real vector $e$, which gives a unique real
$h_b$. Every marked monodromy is real and fixes $e$, so the
same equality defines one function on the base. The construction
is real analytic on each root chart; it is invariant under its
finite deck action and hence descends.

For its coordinate expression, let $H'=H\Gamma_6$, $H(1)=I$.
Take the positive zero meridian with stem from $1$ to $1/32$
and a counterclockwise circle of radius $1/32$, and the infinity
meridian $\lambda=\exp(-2\pi it)$, $0\le t\le1$.
Let $H_0,H_\infty^{\rm mon}$ be their based continuations.
The comparison matrix $X$ is specified by
\[
 XH_0=BX,\qquad XH_\infty^{\rm mon}=AX,\qquad Xe_0=e_0.
\]
Its existence is the marked flat-bundle identification in
\cref{night:P1.15}. The equations $RA=AR$, $RB=BR$ for the
explicit integral matrices force $R$ to be scalar. Thus any
two such comparisons differ by a scalar, and $Xe_0=e_0$ fixes
that scalar. This proves uniqueness with the based meridians
specified.

Put $Y=XH$ and $k=625\pi/96$. Since
$\alpha(w_0)=ki$ and $\alpha(w_1)=\alpha(w_2)=0$,
the vector $kn$ has coordinates
$(0,v_2,v_1,-i,0,0)^t$ in the quotient lift for a unique
pair $(v_2,v_1)$. Equivalently, the lower five coordinates
of its image under $Y$ are real. Equation
\eqref{maps:eq:F.20} gives
\[
 h_b(\lambda)=\frac{96}{625\pi}
 \operatorname{Im}\bigl(Y_{01}v_2+Y_{02}v_1-iY_{03}\bigr).
\]
This proves the existence, uniqueness and normalization of the
real system used in the height computation.
\end{proof}

\begin{remark}[Change of splitting]
\label{rem:V:splitting-change}
Every holomorphic splitting of the pullback extension over $P$
has the form $\sigma_b-a\alpha e$, $a\in\mathbb C$, because
$H^0(P^\vee)=\mathbb C\alpha$. Its parameter and height are
related to the ones fixed here by the exact equalities
\[
 (\sigma_b-a\alpha e+\xi_{\rm alt}\alpha e)(P)
  =(\sigma_b+(\xi_{\rm alt}-a)\alpha e)(P),
 \qquad h_{\rm alt}=h_b-\operatorname{Im}a.
\]
Thus $\xi=\xi_{\rm alt}-a$. We use $\sigma_b$ and $\xi$
throughout.
\end{remark}
